%% Hand-authored circuitikz figure -- NOT generated by maestro2tex.
%% Transcribed from the castalia library cells randles_cell and randles_electrode,
%% read from OA (dbOpenCellViewByType -> instances, instTerms, r/c properties).
%% Component values are in the companion table in note_models.tex; keep the two in
%% step if either cell is edited.
\begin{figure}[htbp]
  \centering
  \resizebox{0.72\linewidth}{!}{%
\ctikzset{bipoles/length=0.95cm, resistors/scale=0.9, capacitors/scale=0.9}
\begin{circuitikz}[
    every node/.append style={font=\footnotesize},
    nlab/.style={font=\scriptsize, inner sep=1.5pt},
    plab/.style={font=\footnotesize\itshape, inner sep=2pt},
    railnode/.style={circle, fill=black, inner sep=1.05pt},
  ]

% ============ (a) three-electrode cell: randles_cell ==================
\def\ya{0}
\node[plab, anchor=north west] at (-0.4,\ya-1.6) {(a) \texttt{randles\_cell}};

\draw (0,\ya) node[railnode]{} node[nlab, above=3pt]{\texttt{CE}}
      to[R, l=$R_{\mathrm{ce}}$] (2.4,\ya);
\draw (2.4,\ya) node[railnode]{} node[nlab, above=3pt]{\texttt{RE}}
      to[R, l=$R_u$] (4.8,\ya);
\node[railnode] at (4.8,\ya){};
\node[nlab, anchor=south east, inner sep=2.5pt] at (4.72,\ya+0.06){$\mathrm{WE_s}$};

% parallel R_ct || C_dl between WE_s and WE
\draw (4.8,\ya) -- (4.8,\ya+0.85);
\draw (4.8,\ya+0.85) to[R, l=$R_{\mathrm{ct}}$] (7.6,\ya+0.85);
\draw (7.6,\ya+0.85) -- (7.6,\ya);
\draw (4.8,\ya) -- (4.8,\ya-0.85);
\draw (4.8,\ya-0.85) to[C, l_=$C_{\mathrm{dl}}$] (7.6,\ya-0.85);
\draw (7.6,\ya-0.85) -- (7.6,\ya);
\draw (7.6,\ya) node[railnode]{} -- (8.5,\ya) node[nlab, right]{\texttt{WE}};

% ============ (b) single electrode: randles_electrode ================
\def\yb{-3.9}
\node[plab, anchor=north west] at (-0.4,\yb-1.6) {(b) \texttt{randles\_electrode}};

\draw (0,\yb) node[railnode]{} node[nlab, above=3pt]{\texttt{a}}
      to[R, l=$R_s$] (2.4,\yb);
\node[railnode] at (2.4,\yb){};

\draw (2.4,\yb) -- (2.4,\yb+0.85);
\draw (2.4,\yb+0.85) to[R, l=$R_{\mathrm{ct}}$] (5.2,\yb+0.85);
\draw (5.2,\yb+0.85) -- (5.2,\yb);
\draw (2.4,\yb) -- (2.4,\yb-0.85);
\draw (2.4,\yb-0.85) to[C, l_=$C_{\mathrm{dl}}$] (5.2,\yb-0.85);
\draw (5.2,\yb-0.85) -- (5.2,\yb);
\draw (5.2,\yb) node[railnode]{} -- (6.1,\yb) node[nlab, right]{\texttt{b}};

\end{circuitikz}
  }
  \caption[{Electrochemical models used by the analog testbenches.}]{Electrochemical models used by the analog testbenches. (a) \texttt{randles\_cell}, the three-electrode cell: solution resistance \(R_{\mathrm{ce}}\) from the counter to the reference electrode, uncompensated resistance \(R_u\) from the reference to the working-electrode surface, and the interface itself as charge-transfer resistance \(R_{\mathrm{ct}}\) in parallel with double-layer capacitance \(C_{\mathrm{dl}}\). (b) \texttt{randles\_electrode}, the same interface as a two-terminal element with its own series resistance \(R_s\). Values are given in Table~\ref{tab:analog-models}. Both are lumped, linear and frequency-independent: no Warburg diffusion term and no potential dependence, so they represent small-signal behaviour about a bias point rather than a full electrochemical response.}
  \label{fig:analog-models}
\end{figure}
